\diarytitle{0099}{$\sin^3(2x^2+1)$ の導関数（連鎖律）}

$u = 2x^{2}+1$ とおくと $y = \sin^{3}u = (\sin u)^{3}$ であり、$u' = 4x$。外側の $(\ \cdot\ )^{3}$、内側の $\sin u$、さらに内側の $u$ という多重合成に連鎖律を繰り返し適用する。
\[
\frac{dy}{du} = 3\sin^{2}u, \qquad \frac{d(\sin u)}{dx} = \cos u\cdot u'
\]
より
\[
y' = 3\sin^{2}u\cdot\cos u\cdot u' = 3\sin^{2}(2x^{2}+1)\cos(2x^{2}+1)\cdot 4x
\]
整理すると
\[
\boxed{\; y' = 12x\sin^{2}\!\left(2x^{2}+1\right)\cos\!\left(2x^{2}+1\right) \;}
\]

（検算: $x=0$ のとき $u=1$ で $y'(0) = 12\cdot 0\cdot(\cdots) = 0$。実際 $y(x) = \sin^{3}(2x^{2}+1)$ は偶関数（$x\to -x$ で不変）なので、微分可能ならば $y'(0)=0$ でなければならず、整合している。）
